\section*{Hoofdstuk \ref{ch:b2:waves-in-media} --- Elektromagnetische golven
in plasma's, geleiders en diëlektrica}

\begin{solution}{exo:b2:waves-in-media:1}
$f_p \approx 9\sqrt n$ Hz: \qty{2.8}{MHz}, \qty{9}{MHz}, \qty{9}{GHz},
\qty{2.6e15}{Hz} ($\lambda_p = \qty{115}{nm}$), \qty{2.8}{kHz}. AM op
\qty{1}{MHz}: alleen door de interstellaire ruimte. FM: door de ionosfeer,
niet door het \omterm{def:b2:waves-in-media:plasma}{plasma} bij terugkeer. Gps: door alles behalve dat \omterm{def:b2:waves-in-media:plasma}{plasma}.
Licht: door alles behalve koper.
\end{solution}

\begin{solution}{exo:b2:waves-in-media:2}
Koper: \qty{9.2}{mm}, \qty{2.1}{mm}, \qty{65}{\micro m}, \qty{2.1}{\micro
m}; aluminium bij \qty{1}{MHz}: \qty{84}{\micro m}; zeewater: \qty{2.3}{m}
bij \qty{10}{kHz}, \qty{7}{mm} bij \qty{1}{GHz}; grond: \qty{5}{m} bij
\qty{1}{MHz}. Een onderzeeër op \qty{10}{m} heeft
$\delta
\gtrsim \qty{10}{m}$ nodig: onder ongeveer \qty{500}{Hz}.
\end{solution}

\begin{solution}{exo:b2:waves-in-media:3}
$n = 1.528$, $1.515$, $1.509$; $v = 1.963$, $1.980$, $\qty{1.987e8}{m/s}$;
$n_g = n +
2B/\lambda^2 = 1.545$. Brekingshoeken $\arcsin(0.707/n)$:
\qty{27.57}{\degree} en \qty{27.93}{\degree}: een spreiding van
\qty{0.36}{\degree}.
\end{solution}

\begin{solution}{exo:b2:waves-in-media:4}
$f_p = \qty{4.0}{MHz}$: \qty{5}{MHz} gaat erdoor; \qty{2}{MHz} wordt
teruggekaatst en dringt $c/\sqrt{\omega_p^2 - \omega^2} = \qty{14}{m}$ door.
Voor $n = \qty{1.5e12}{m^{-3}}$: \qty{11}{MHz}. De dichtheid volgt de
ioniserende straling van de zon --- uur, seizoen en de elfjarige cyclus ---
zodat de hoogst bruikbare frequentie daarmee meebeweegt.
\end{solution}

\begin{solution}{exo:b2:waves-in-media:5}
(a) $1/\gamma\pi a^2 = \qty{5.3}{m\ohm/m}$. (b)
$\delta = \qty{65}{\micro m}$;
$R \approx 1/\gamma2\pi a\delta
= \qty{41}{m\ohm/m}$, $7.7$ keer de
gelijkstroomwaarde. (c) $\delta = \qty{6.5}{\micro m}$, \qty{0.41}{\ohm/m},
$77$ keer. (d) $R \propto \sqrt f$ groeit even snel als $L\omega$ zodra de
huid overheerst, en elke winding voegt weerstand toe; litzedraad bundelt
vele geïsoleerde aders die dunner zijn dan $\delta$, zodat al het koper
geleidt.
\end{solution}

\begin{solution}{exo:b2:waves-in-media:6}
$\delta_{\text{Al}} = 11.9$, $0.84$, \qty{0.084}{mm}: $d/\delta = 0.08$,
$1.2$, $12$: $0.7$, $10$, \qty{103}{dB}. Bij \qty{50}{Hz} is de geleider
doorzichtig; ijzer met hoge $\mu_r$ leidt de flux om (magnetische
afscherming) en verkleint $\delta$ bovendien met $\sqrt{\mu_r}$.
\end{solution}

\begin{solution}{exo:b2:waves-in-media:7}
(a) $v_0 = eE_0/m\omega$;
$\tfrac12nm\langle v^2\rangle = ne^2E_0^2/4m\omega^2 = \tfrac14\varepsilon_0
E_0^2\,\omega_p^2/\omega^2$.
(b) $\tfrac14\varepsilon_0E_0^2$ en
$\tfrac14\varepsilon_0E_0^2(1 - \omega_p^2/
\omega^2)$. (c) De som is
$\tfrac12\varepsilon_0E_0^2$ (de haakjes geven $1$);
$\langle\Pi
\rangle = E_0^2k/2\mu_0\omega$, en de verhouding is
$c^2k/\omega = v_g$. (d) $\tfrac14\cdot\tfrac14
/\tfrac12 = 1/8$.
\end{solution}

\begin{solution}{exo:b2:waves-in-media:8}
(a)
$1/v_g = (1/c)(1 - \omega_p^2/\omega^2)^{-1/2} \approx (1 + \omega_p^2/2\omega^2)/c$.
(b) $\Delta t
= (L\omega_p^2/2c)(1/\omega_1^2 - 1/\omega_2^2)$ met
$\omega_p^2 = ne^2/\varepsilon_0m$ en $\omega =
2\pi f$: constante
$e^2/8\pi^2\varepsilon_0mc = \qty{1.34e-7}{}$ (SI). (c)
$nL = \qty{9.3e23}{m^{-2}}$, $(1/f_1^2 - 1/f_2^2) = \qty{5.7e-18}{s^2}$:
$\Delta t = \qty{0.72}{s}$. (d) De kolom $nL$ (de dispersiemaat), en daarmee
de afstand bij een model voor $n$.
\end{solution}

\begin{solution}{exo:b2:waves-in-media:9}
(a)
$\operatorname{Im}\underline\chi = \omega_p^2\Gamma\omega/[(\omega_0^2 - \omega^2)^2 + \Gamma^2\omega^2]$
met $\omega_0^2 - \omega^2 \approx 2\omega_0(\omega_0 - \omega)$:
$n'' = \tfrac12\operatorname{Im}\underline\chi$ is het genoemde
lorentzprofiel. (b) $n''(\omega_0) = \omega_p^2/2\omega_0\Gamma$,
$\alpha = 2n''\omega_0/c =
\omega_p^2/\Gamma c$. (c)
$\omega_p^2 = \qty{3.2e20}{s^{-2}}$: $\alpha = \qty{1.8e4}{m^{-1}}$;
$\ln10/\alpha = \qty{0.13}{mm}$. (d)
$n' - 1 \propto (\omega_0 - \omega)/[(\omega_0 - \omega)^2 +
\Gamma^2/4]$:
positief eronder, negatief erboven, en dalend binnen $\pm\Gamma/2$ van de
lijn --- daar is zij anomaal.
\end{solution}

\begin{solution}{exo:b2:waves-in-media:10}
(a) $\underline E = E_0\eu^{-x/\delta}\eu^{\iu(\omega t - x/\delta)}$,
$\underline B = \underline k\underline E/\omega
= (1 - \iu)\underline E/\omega\delta$.
(b)
$\langle\Pi_x\rangle = \tfrac12\operatorname{Re}(\underline E
\underline B^*)/\mu_0 = E_0^2/2\mu_0\omega\delta$
in $x = 0$. (c)
$\int_0^\infty\tfrac12\gamma E_0^2
\eu^{-2x/\delta}\dd x = \gamma E_0^2\delta/4 = E_0^2/2\mu_0\omega\delta$
met $\gamma = 2/\mu_0\omega\delta^2$. (d)
$H_0 = \sqrt2E_0/\mu_0\omega\delta$, zodat
$\tfrac12R_sH_0^2 = E_0^2/\gamma\mu_0^2\omega^2\delta^3 =
E_0^2/2\mu_0\omega\delta$.
Koper bij \qty{1}{GHz}: $R_s = 1/\gamma\delta = \qty{8}{m\ohm}$;
$\tfrac12 \times 8 \times 10^{-3} \times 10^4 = \qty{40}{W/m^2}$.
\end{solution}

\begin{solution}{exo:b2:waves-in-media:11}
(a) $\omega_p = \qty{1.4e16}{rad/s}$,
$\lambda_p = 2\pi c/\omega_p = \qty{140}{nm}$. (b)
$\omega = \qty{3.8e15}{rad/s}
< \omega_p$:
$1/\kappa = c/\sqrt{\omega_p^2 - \omega^2} = \qty{23}{nm}$, totale
reflectie. (c) Nog altijd onder $\omega_p$: spiegelend; doorzichtigheid pas
voorbij \qty{140}{nm} (bij natrium \qty{210}{nm}). (d) Botsingen geven
$\underline\gamma$ een reëel deel en enkele procenten absorptie; resonanties
van gebonden elektronen in het zichtbare (goud, koper) slikken het blauw en
kleuren het metaal.
\end{solution}

\begin{solution}{exo:b2:waves-in-media:12}
(a) $\omega\tau = 0.12$: $\underline\varepsilon_r = 84 - 9.7\iu$; bij
\qty{20}{GHz} is $\omega\tau = 1$: $45 - 40\iu$. (b)
$\underline n \approx 9.2 - 0.53\iu$; $c/2n''\omega = \qty{1.8}{cm}$: de
oven kookt de buitenste twee centimeter, de rest doet de geleiding. (c)
$\tfrac12 \times
1.54 \times 10^{10} \times 8.85 \times 10^{-12} \times 9.7 \times 4 \times 10^6 = \qty{2.6}{MW/m^3}$.
(d) Bij \qty{20}{GHz} zou de absorptielengte een millimeter zijn: alleen de
buitenlaag zou koken; \qty{2.45}{GHz} dringt door, en is een vrije band.
\end{solution}

\begin{solution}{pb:b2:waves-in-media:1}
\textbf{1.} $f_p = \qty{9}{MHz}$: gps komt erdoor, de kortegolf wordt teruggekaatst.

\textbf{2.} $k = (\omega/c)\sqrt{1 - \omega_p^2/\omega^2}$:
$v_\varphi = \omega/k \approx c(1 + \omega_p^2/
2\omega^2)$,
$v_g = \dd\omega/\dd k = c^2k/\omega \approx c(1 - \omega_p^2/2\omega^2)$.

\textbf{3.}
$\Delta t = \int(1/v_g - 1/c)\dd z = \int\omega_p^2\dd z/2c\omega^2 = e^2N_T/2\varepsilon_0
mc\omega^2 = (e^2/8\pi^2\varepsilon_0mc)N_T/f^2 = 1.34 \times 10^{-7} \times 2 \times 10^{17}/2.48 \times
10^{18} = \qty{11}{ns}$:
\qty{3.2}{m}.

\textbf{4.} $N_T = (\Delta t_1 - \Delta t_2)/K(1/f_1^2 - 1/f_2^2)$; het
verschil is \qty{7}{ns}: voor $1\%$ is \qty{0.07}{ns} nodig.

\textbf{5.} $v_\varphi$ overtreft $c$ met evenveel als $v_g$ tekortkomt: de
fase van de draaggolf komt te vroeg aan terwijl de modulatie te laat komt.
Een ontvanger die de fase van de draaggolf gebruikt, past de correctie met
het omgekeerde teken toe.

\textbf{6.} \qty{3.2}{m} overdag, \qty{0.3}{m} 's nachts.

\textbf{7.} $f_p = \qty{0.9}{MHz}$ met botsingen: midden- en kortegolf
worden geabsorbeerd (een radiostilte); gps op \qty{1.5}{GHz} ondervindt een
absorptie $\propto 1/\omega^2$, verwaarloosbaar.

\textbf{8.} $f_1\Delta t = 1.575 \times 10^9 \times 1.1 \times 10^{-8} = 17$
perioden.

\textbf{9.} $\delta = \sqrt{2/\mu_0\gamma\omega} = \qty{10}{\micro m}$: een
wand van een millimeter is honderd \omterm{def:b2:dispersion-wave-packets:complex}{indringdiepten}.

\textbf{10.} $R_s = 1/\gamma\delta = \qty{0.1}{\ohm}$;
$\tfrac12R_sH_0^2 = \qty{45}{W/m^2}$: \qty{45}{W}, $6\%$.

\textbf{11.} $\eu^{-\pi \times 0.5} = 0.21$: \qty{-14}{dB} in amplitude,
\qty{-27}{dB} in vermogen (en het kleine gat koppelt sowieso al slecht).

\textbf{12.} Licht, $\lambda = \qty{0.5}{\micro m}$, is tweeduizend keer
kleiner dan de gaten en gaat erdoor; de microgolf, $\lambda = \qty{12}{cm}$,
is honderd keer groter.

\textbf{13.} $\underline n = 2 - 0.005\iu$; $c/2n''\omega = \qty{1.9}{m}$:
het glas wordt nauwelijks warm.

\textbf{14.} $j \approx H_0/\delta = \qty{3e7}{A/m^2}$,
$j^2/\gamma = \qty{9e8}{W/m^3}$ in een laag van tien micrometer: de punten
worden binnen milliseconden heet, zenden elektronen uit en ioniseren de
lucht --- vonken.

\textbf{15.} $\delta = \qty{18}{\micro m}$; doorsnede
$2\pi a\delta = \qty{5.6e-8}{m^2}$ tegen $\pi a^2 = \qty{7.9e-7}{m^2}$:
$R_{\text{ac}} = \qty{0.30}{\ohm/m}$, veertien keer de gelijkstroomwaarde.

\textbf{16.} $RI^2 = \qty{120}{W/m}$; oppervlak $\pi d = \qty{3.1e-3}{m^2}$
per meter: $\Delta T \approx \qty{1900}{K}$ --- hij zou smelten; zulke
spoelen zijn watergekoelde buizen.

\textbf{17.} $\delta = \qty{0.14}{mm}$: de opgewekte stroom, en de warmte,
blijven in een tiende millimeter; een korte puls gevolgd door afschrikken
hardt het oppervlak van een tandwiel terwijl de kern taai blijft.

\textbf{18.} $R = 1/\gamma2\pi a\delta = \qty{0.10}{\ohm/m}$, drie keer
minder; de binnenkant voert niets --- vandaar een buis, met het koelwater
waar het koper toch verspild zou zijn.

\textbf{19.} $\delta = a$ bij $f = 1/\pi\mu_0\gamma a^2 = \qty{17}{Hz}$;
daaronder is de stroom uniform en geldt de gelijkstroomformule.

\textbf{20.}
$n^2 \approx 1 + \omega_p^2/\omega_0^2 + (\omega_p^2/\omega_0^2)(\lambda_0/\lambda)^2$:
de vorm van Cauchy met $n(\infty)^2 - 1 = \omega_p^2/\omega_0^2 = 1.25$ voor
$n(\infty) = 1.5$.

\textbf{21.} $n = 1.531$ en $1.510$, $v = 1.96$ en \qty{1.99e8}{m/s};
$n_g(550) =
1.550$; $n_g(400) - n_g(700) = 0.063$: \qty{210}{ns} over een
kilometer.

\textbf{22.} $\sin((A + D)/2) = 1.5165 \times 0.5$:
$D = \qty{38.6}{\degree}$;
$\dd D = 2\sin
(A/2)\dd n/\cos((A + D)/2) = 1.53\,\dd n$:
$\Delta D = \qty{1.85}{\degree}$, \qty{6.5}{cm} op \qty{2}{m}.

\textbf{23.} Bij \qty{200}{nm} nadert de frequentie de resonantie en groeit
$n''$: ondoorzichtig. Een zwakke resonantie bij \qty{700}{nm} slikt het
rood: groen glas.

\textbf{24.} Blauw ligt dichter bij de ultraviolette resonantie: de gebonden
elektronen reageren sterker, $n$ is groter, en de afwijking dus ook.

\textbf{25.} F-laag bij \qty{1.5}{GHz}: imaginaire $\underline\gamma$,
voortplanting met een kleine groepsvertraging. Staal bij \qty{2.45}{GHz}:
reële $\underline\gamma$, \omterm{thm:b2:waves-in-media:skin}{huideffect}, reflectie. Koper bij \qty{13.56}{MHz}:
hetzelfde, met een huid van \qty{18}{\micro m}. Glas in het zichtbare:
$\iu\omega\varepsilon_0\chi$, een reële brekingsindex, een doorzichtig
\omterm{def:b2:dispersion-wave-packets:relation}{dispersief medium}.
\end{solution}
